\chapter{Высота и ожидание}\label{ch:vysota}

\textbf{Что вы узнете:} как менять высоту и как заставить дрон ждать.

\section{Команда set\_z}

\kod{d.set_z(z, seconds)} переводит дрона на высоту \kod{z} сантиметров за \kod{seconds} секунд. Дрон при этом остаётся над той же клеткой.

\solfile{l04_height.py}

Что происходит: взлёт на 60\cm, подъём до 120\cm (3\,с), потом ещё 2\,с на той же высоте, спуск до 60\cm (3\,с) и посадка.

\section{Ожидание}

Отдельной команды «ждать» нет. Если надо подождать, дрону говорят: «перейди на ту высоту, на которой ты уже висишь». Вторая строка \kod{d.set_z(120, 2)} в примере выше и есть ожидание: высота не меняется, проходят две секунды.

\begin{zametka}
Ожидание = \kod{set_z} на той же высоте. Это основной приём для шоу из нескольких дронов (урок~\ref{ch:ochered}).
\end{zametka}

\section{Что нельзя}

\begin{itemize}
  \item Высота не больше лимита (по умолчанию 150\cm, максимум 200\cm).
  \item Одно ожидание не длиннее 60 секунд. Если нужно дольше, разбейте на несколько шагов. Эта проверка срабатывает при загрузке в дрон.
\end{itemize}

\begin{probuy}
  \item Лесенка вверх: взлёт на 50\cm, потом ступени 100\cm и 150\cm, по 2 секунды каждая.
  \item «Дыхание»: между 60\cm и 120\cm два раза туда и обратно, по 2 секунды на каждый переход.
  \item Дрон висит 10 секунд на 100\cm: первые пять секунд светится синим, следующие пять красным.
  \item «Лифт»: взлёт на 150\cm, спуск до 30\cm за 4 секунды, потом подъём до 100\cm за 3 секунды.
\end{probuy}
\otvety{4}
